\input header
\input ras
\def\GO{{G\!O}}
\def\TF{{\cal E}}
\def\T{{\bf T}}

\footline={\hss\tenrm 3-\folio\hss}\pageno=1

\centerline{\bf THE JULIA SET}\medskip

\centerline{\bf \S3. Fatou and Julia: Dynamics on the Riemann Sphere.}\bigskip

The local study of iterated holomorphic mappings, in a neighborhood of
a fixed point, was quite well developed in the late $19^{th}$ century.
(Compare \S\S6,7.) However, except for one very simple case studied
by Cayley, nothing was known about the global behavior of iterated
holomorphic maps
until 1906, when Pierre Fatou described the following startling example.
For the map \[z\mapsto z^2/(z^2+2)\], he showed that
almost every orbit under iteration
converges to zero, even though there is a Cantor set of exceptional points
for which the orbit remains bounded away from zero. (Problem 3-6.)
This aroused great interest. After a hiatus dug
the first world war, the subject was taken up in depth by Fatou, and also
by Gaston Julia and others such as S. Latt\'es and J. F. Ritt. The most
fundamental and incisive contributions were those of Fatou himself,
although Julia developed much closely related material at more or less the
same time. Julia, who had been badly wounded
during the war, was awarded the ``Grand Prix des Sciences math\'ematiques''
by the Paris Academy of Sciences in 1918 for his work.\medskip

{\bf Definition.} Let \[S\] be a Riemann surface, let \[f:S\to S\]
be a non-constant holomorphic mapping, and let \[f^{\circ n}:S\to S\] be its
\[n$-fold iterate.
Fixing some point \[z_0\in S\], we have the following basic dichotomy: {\it
If there exists some neighborhood \[U\] of \[z_0\] so that the sequence of
iterates \[\{f^{\circ n}\}\] restricted to \[U\] forms a normal family, then
we say that \[z_0\] is a {\bit regular\/} or {\bit normal\/} point, or that
\[z_0\] belongs to the {\bit Fatou set\/} of \[f\]. Otherwise, if no such
neighborhood exists, we say that \[z_0\] belongs
to the {\bit Julia set\/} \[J=J(f)\].} (For sharper formulations of this
dichotomy in the rational case, see \S11.8 and Problem 13-1.)\smallskip

Thus, by its very definition, the Julia set \[J\] is a closed subset of
\[S\], while the Fatou set \[S\ssm J\] is the complementary open
subset. (The choice
as to which of these two sets should be named after Julia and which after
Fatou is rather arbitrary. The term ``Julia set'' is firmly established,
but the Fatou set is often called by other names, such as ``stable set''
or ``normal set''.) {\it
Roughly speaking,
\[z_0\] belongs to the Fatou set if dynamics in some neighborhood of \[z_0\]
is relatively tame, and belongs to the Julia set, if dynamics in every
neighborhood of \[z_0\] is more wild.}\smallskip

The classical example, and the one which we will emphasize, is the
case where \[S\] is the Riemann sphere \[\hat\C=\C\cup\infty\].
Any holomorphic map \[f:\hat\C\to \hat\C\] on the Riemann sphere
can be expressed as a rational
function, that is as the quotient \[f(z)=p(z)/q(z)\] of two polynomials.
Here we may assume that \[p(z)\] and \[q(z)\] have no common roots.
The {\bit degree} \[d\] of \[f=p/q\] is then equal to the maximum of the
degrees of \[p\] and \[q\]. In particular, for almost every choice of constant
\[c\in\hat\C\] the equation \[f(z)=c\] has exactly \[d\] distinct solutions
in \[\hat\C\]. (For every choice of \[c\] it has at
least one solution, since we assume that \[d>0\].)\smallskip

As a simple example, consider the squaring map \[s:z\mapsto z^2\] on
\[\hat\C\]. The entire open disk \[D\] is contained in the Fatou set of
\[s\], since successive iterates on any compact subset converge uniformly
to zero. Similarly the exterior \[\hat\C\ssm \bar D\] is contained in the
Fatou set, since the iterates of \[s\] converge to the constant function
\[z\mapsto\infty\] outside of \[\bar D\]. On the other hand, if \[z_0\]
belongs to the unit circle, then in any neighborhood of \[z_0\]
any limit of iterates \[s^{\circ n}\] would necessarily
have a jump discontinuity as we cross the unit circle.
This shows that the Julia set \[J(s)\] is precisely equal to the unit circle.
\smallskip

Such smooth Julia sets are rather exceptional. (Compare \S5.)
See Figure 1 for some much more typical examples of Julia sets for polynomial
mappings. Figure 1a
shows a rather wild Jordan curve,
Figure 1b a rather thick Cantor set, Figure 1c a ``dendrite'', 
and Figure 1d a more complicated example, the ``airplane'', with a
superattracting period 3 orbit. (Further examples of Julia sets are shown
in Figures 2-5, 8-10, 12, and 17.)

\pageinsert % \eject

\vfil
\insertRaster qj_quasicircle.ps pixels 720 by 900 scaled 320
\smallskip

\centerline{Figure 1a. A simple closed curve,}
\centerline{Julia set for \[z\mapsto z^2+(.99+.14i)z\].}
\vfil
\insertRaster qj_cantor.ps pixels 960 by 540 scaled 320
\smallskip
\centerline{Figure 1b. A totally disconnected Julia set,}
\centerline{\[z\mapsto z^2+(-.765+.12i)\].}
\vfil
\endinsert % \eject

\pageinsert
\vfil
\insertRaster qj_i.ps pixels 960 by 960 scaled 400
\smallskip

\centerline{Figure 1c. A ``dendrite'', Julia set for \[z\mapsto z^2+i\].}
\vfil%\endinsert % \eject


\insertRaster qj_airplane.ps pixels 960 by 240 scaled 400
\smallskip
\centerline{Figure 1d. Julia set for \[z\mapsto z^2-1.75488\ldots\],
the ``airplane''.}
\vfil\endinsert % \eject

\smallskip

We will also need the following concept.\smallskip

{\bf Definition.} By the {\bit grand orbit} of a point
\[z\] under \[f:S\to S\] we mean the
set \[\GO(z,f)\] consisting of all points \[z'\in S\] whose orbits
eventually intersect the orbit of \[z\]. Thus \[z\] and \[z'\] have the same
grand orbit if and only if \[f^{\circ m}(z)=f^{\circ n}(z')\] for some
choice of \[m\ge 0\] and \[n\ge 0\].\smallskip

Here are some basic properties of the Julia set.

{\QP{\bf 3.1. Lemma.} {\it The Julia set \[J(f)\] of a
holomorphic map \[f:S\to S\] is {\bit fully invariant}
under \[f\]. That is,
if \[z\] belongs to \[J(f)\], then the entire grand orbit \[\GO(z,f)\]
is contained in \[J(f)\].}\medskip}

Evidently it suffices to prove that \[z\in J(f)\] if and only
if \[f(z)\in J(f)\].
A completely equivalent statement is that the Fatou set is fully invariant.
The proof, making use of the fact that a non-constant holomorphic map
takes open sets to open sets, is completely straightforward
and will be left to the reader.\QED

It follows that the Julia set possesses a great deal of {\bit
self-similarity\/}: {\it Whenever\break
\[f(z_1)=z_2\] in \[J(f)\], with derivative \[f'(z_1)\ne 0\], there
is an induced conformal\break isomorphism from a neighborhood \[N_1\] of \[z_1\]
to a neighborhood \[N_2\] of \[z_2\] which takes \[N_1\cap J(f)\] precisely
onto \[N_2\cap J(f)\].} (Compare Problem 3-7.)

{\QP {\bf 3.2. Lemma.} {\it For any \[n>0\], the Julia set \[J(f^{\circ n})
\] of the \[n$-fold iterate coincides with the Julia set \[J(f)\].}\medskip}

\noindent Again, the proof will be left to the reader.\QED

Now consider a {\bit periodic orbit\/} or ``{\bit cycle\/}''
$$	f\;:\;z_0\mapsto z_1\mapsto\;\cdots\;\mapsto z_{n-1}
	\mapsto z_n=z_0\,. $$
If the points \[z_1\,,\,\ldots\,,\,z_n\]
are all distinct, then the integer \[n\ge 1\] is called the {\bit period\/}.
If the Riemann surface \[S\] is \[\C\] (or an open subset
of \[\C\]), then the derivative
$$	\lambda\= (f^{\circ n})'(z_i)\=f'(z_1)\cdot f'(z_2)\cdots
	 f'(z_n) $$
is a well defined complex number called the {\bit multiplier~} or the
{\bit eigenvalue~}  of this periodic orbit. More generally, for
self-maps of an arbitrary Riemann surface the multiplier of a periodic
orbit can be defined using a local coordinate chart around any point of the
orbit. By definition, a
periodic orbit is either {\bit attracting~} or {\bit repelling~}
or {\bit indifferent\/} (\[=\]{\bit neutral~}) according as its multiplier
satisfies \[|\lambda|<1\] or \[|\lambda|>1\] or \[|\lambda|=1\]. The orbit
is called {\bit superattracting~} if \[\lambda=0\].

{\bf Caution:} In the special case where the point at infinity is periodic
under a rational map,
\[f^{\circ n}(\infty)=\infty\], this definition may be confusing.
The multiplier \[\lambda\] is {\it
not\/} equal to the limit as \[z\to\infty\] of the derivative of \[f^{\circ n}
(z)\], but is rather equal to
the reciprocal of this number. In fact if we introduce
the local uniformizing parameter \[w=1/z\] for \[z\] near \[\infty\],
then it is easy to check that the derivative of \[w\mapsto 1/f(1/w)\]
at \[w=0\] is
equal to \[\lim_{z\to\infty} 1/(f^{\circ n})'(z)\]. As an
example, if \[f(z)=2z\] then
\[\infty\] is an attracting fixed point with multiplier \[\lambda=1/2\].

In the
case of an attracting periodic orbit, we can define the {\bit basin of
attraction~} to be the open set \[\Omega\subset S\] consisting of all
points \[z\in S\] for which the successive iterates \[f^{\circ n}(z)\,,\,
f^{\circ 2n}(z)\,,\,\ldots\] converge towards some point of the periodic
orbit. In particular, this basin of attraction is defined
in the superattracting case. \smallskip

{\QP{\bf 3.3. Theorem.} {\it  Every attracting periodic orbit is contained
in the Fatou set. In fact the entire basin of
attraction \[\Omega\] for an attracting periodic orbit is contained in the
Fatou set. However the {\bit boundary} \[\partial \Omega\] is contained in the
Julia set, and every repelling periodic orbit is contained in the Julia set.}
\medskip}

{\bf Proof.} In view of 3.2, we need only consider the case of a fixed
point \[f(z_0)=z_0\]. If \[z_0\] is attracting, then it follows
from Taylor's Theorem that the successive iterates of
\[f\], restricted to a small neighborhood of \[z_0\], converge uniformly
to the constant function \[z\mapsto z_0\]. The corresponding statement for
any compact subset of the basin \[\Omega\] then follows easily.
On the other hand, around a
boundary point of this basin, that is a point which belongs to
the closure \[\bar \Omega\] but not to \[\Omega\] itself, it is clear that no
sequence of iterates can converge to a continuous limit. (See Problem
3-2 for a sharper statement.) If \[z_0\] is repelling, then
no sequence of iterates can converge uniformly near \[z_0\], since the
derivative \[df^{\circ n}(z)/dz\] at \[z_0\] takes the value \[\lambda^n\],
which diverges to infinity as \[n\to\infty\]. (Compare \S1.4.)\QED

The case of an indifferent periodic point is much more difficult. (Compare
\S\S7, 8.) One particularly important case is the following.\smallskip

{\bf Definition.} A periodic point \[f^{\circ n}(z_0)=z_0\] is called
{\bit parabolic~} if the multiplier \[\lambda\] at \[z_0\] is equal to
\[+1\], yet \[f^{\circ n}\] is not the identity map, or more generally
if \[\lambda\] is a root of unity, yet no iterate of \[f\] is the identity
map.

As an example, the two fixed points of the rational map \[f(z)=z/(z-1)\]
both have multiplier equal to \[-1\]. These do not count as parabolic
points since \[f\circ f(z)\] is identically equal to \[z\]. This exclusion
is necessary so that the following assertion will be true.

{\QP{\bf 3.4. Lemma.} \it Every parabolic periodic point belongs to the Julia
set.\medskip}

{\bf Proof.} Let \[w\] be a local uniformizing parameter, with \[w=0\]
corresponding to the periodic point. Then some iterate \[f^{\circ m}\]
corresponds to a local mapping of the \[w$-plane with
power series expansion of the form \[\;w\,\mapsto\, w+a_kw^k+a_{k+1}w^{k+1}+
\cdots\,\], where \[k\ge 2\,,\;\; a_k\ne 0\]. It follows that \[f^{\circ mp}\]
corresponds to a power series \[w\mapsto w+pa_kw^k+\cdots\]. Thus the \[k$-th
derivatives of \[f^{\circ mp}\] diverge as \[p\to\infty\].
It follows from 1.4 that
no subsequence can converge locally uniformly.\QED

Now and for the rest of \S3,
let us specialize to the case of a rational map \[f:\hat\C\to\hat\C\] of
degree \[d\ge 2\].

{\QP{\bf 3.5. Lemma.} {\it If \[f\] is rational of degree two or more, then the
Julia set \[J(f)\] is non-vacuous.}\medskip}

{\bf Proof.} If \[J(f)\] were vacuous, then some sequence of iterates
\[f^{\circ n(j)}\] would converge, uniformly over the entire sphere \[\bar
\C\], to a holomorphic limit \[g:\hat\C\to\hat\C\]. (Here we are using the
fact that normality is a local property: Problem 2-5.) A standard topological
argument would then show that the degree of \[f^{\circ n(j)}\] is equal to the
degree of \[g\] for large \[j\]. But the degree of \[f^{\circ n}\] is equal to
\[d^n\], which diverges to infinity as \[n\to\infty\].\QED

A different, more constructive proof of this Lemma will be given in \S9.5.
\smallskip

{\bf Definition.} A point \[z\in \hat\C\] is called
{\bit grand orbit finite} or (to use the classical terminology)
{\bit exceptional}$\,$ under \[f\] if its grand orbit \[\GO(z,f)\subset\hat
\C\] is a finite set. Using Montel's Theorem, we prove the following.

{\QP{\bf 3.6. Lemma.} {\it If \[f\] is rational of
degree two or more, then the set
\[\TF(f)\] of grand orbit finite points can have at most two elements. These
grand orbit finite points,
if they exist, must always be critical points of \[f\], and must belong
to the Fatou set.}\medskip}

{\bf Proof.} (Compare Problem 3-3.)
Note that \[f\] maps any grand orbit \[\GO(z,f)\] onto
itself. Hence, any finite grand orbit must constitute a single
periodic orbit under \[f\]. Each point \[z\] in this finite orbit must
be critical (and in fact
\[(d-1)$-fold critical, where \[d\] is the degree), since otherwise
\[f(z)\] would have two or more pre-images. Therefore, such an
orbit must be attracting, and hence contained in the Fatou set.

If there were three distinct grand orbit finite points, then the union of the
grand orbits of these points would form a finite set whose complement
\[U\] in \[\hat\C\] would be Hyperbolic, with \[f(U)= U\]. Therefore,
the set of iterates of \[f\] restricted to \[U\] would be normal by Montel's
Theorem. Thus both \[U\] and its complement would be contained in the
Fatou set, contradicting 3.5.\QED

{\QP{\bf 3.7. Lemma.} {\it Let \[z_1\] be any point of the Julia set
\[J(f)\subset\hat\C\]. If \[N\]
is a sufficiently small neighborhood of \[z_1\], then the union \[U\] of the
forward images \[f^{\circ n}(N)\] is precisely equal to the complement
\[\hat\C\ssm\TF(f)\] of the set of grand orbit finite points.}\medskip}

In particular, this union \[U\] contains the Julia set \[J(f)\].
(In \S11.2 we will see that just one forward image \[f^{\circ n}(N)\] actually
contains the entire Julia set, provided that \[n\] is sufficiently large.)

{\bf Proof of 3.7.} Let \[E\] be the complementary set \[\hat\C\ssm U\].
We have \[f(U)\subset U\], or equivalently \[f^{-1}(E)\subset E\], by the
construction. Since \[U\] intersects the Julia set, it follows from Montel's
Theorem that its complement \[E\] has at most two points. Now since \[E\] is
finite and \[f\] is onto, a counting argument shows that \[f^{-1}(E)=E\],
hence \[E\] is contained in the set \[\TF(f)\] of grand orbit
finite points. If the initial neighborhood \[N\]
is small enough to be disjoint from \[\TF(f)\], then it follows that
\[E=\TF(f)\].\QED

{\QP{\bf 3.8. Corollary.} {\it If the Julia set contains an interior point
\[z_1\], then it must be equal to the entire Riemann sphere.}\medskip}

For if we choose a neighborhood \[N\subset J\], then 3.7 shows that the
union \[U\subset J\] of forward images of \[N\] is everywhere dense on
\[\hat\C\].\QED

{\QP{\bf 3.9. Theorem.} {\it
If \[z_0\in J(f)\], then the set of all iterated pre-images
$$	\{z\;:\; f^{\circ n}(z)=z_0\quad{\rm for\;some\quad} n\ge0\}	$$
is everywhere dense in \[J(f)\]. In particular, it follows that the
grand orbit\break \[\GO
(z_0\,,f)\] is everywhere dense in \[J(f)\].}\medskip}

For \[z_0\] is not a grand orbit finite point, so 3.7 shows that every point
\[z_1\in J(f)\] can be approximated arbitrarily closely by points \[z\]
whose forward orbits contain \[z_0\].\QED

This Theorem suggests an algorithm for computing pictures of the
Julia set: Starting with any \[z_0\in J(f)\], first compute all pre-images
\[f(z_1)=z_0\], then compute all pre-images \[f(z_2)=z_1\], and so on, thus
eventually coming arbitrarily close to every point of \[J(f)\]. This method
is most often used in the quadratic case, since quadratic equations are very
easy to solve. The method is very insensitive to round-off errors; since
\[f\] tends to be expanding on its Julia set, so that \[f^{-1}\] tends to be
contracting. (Compare Problem 3-4. ??)

{\QP {\bf 3.10. Corollary. } {\it If \[f\] has degree two or more, then
\[J(f)\] has no isolated points.}\medskip}

{\bf Proof.} Since \[J(f)\] is fully invariant, it follows from 3.5 and 3.6
that \[J(f)\] must be an infinite set. Hence it contains at least one
limit point \[z_0\]. Now the iterated pre-images of \[z_0\] form a dense set
of non-isolated points in \[J(f)\].\QED

A property of a point of \[J\] is said to be true for {\bit generic}
\[z\in J\] if it is true for all points in some countable intersection
of dense open subsets \[U_i\cap J\subset J\]. (Compare \S8. Here the notation
is supposed to indicate that the \[U_i\] are open subsets of \[\hat\C\]
and that the closure of \[U_i\cap J\] is equal to \[J\].) By Baire's
Theorem, any such countable intersection is itself a dense subset of \[J\].

{\QP{\bf 3.11. Corollary.} \it For generic \[z\in J(f)\], the forward orbit
\hfil\break \[\{z\,,\,f(z)\,,\,f^{\circ 2}(z)\,,\,\ldots\;\}\;\] is everywhere
dense in \[J(f)\].\medskip}

{\bf Proof.} Let \[\{B_j\}\] be a countable collection of open sets
forming a basis for the topology of \[\hat\C\]. For each \[B_j\] which
intersects \[J=J(f)\], let \[U_j\] be the union of the iterated pre-images
\[f^{-n}(B_j)\] for \[n\ge 0\]. Then it follows from 3.9 that the closure
of \[U_j\cap J\] is equal to the entire Julia set \[J\], and the conclusion
follows.\QED
\pageinsert
\vfil
\insertRaster cubicj_example.ps pixels 960 by 1120 scaled 400
\bigskip\bigskip
\centerline{Figure 2. Julia set for \[f(z)=z^3+{12\over 25}z+{116\over 125}i\].}
\vfil
\endinsert

We will continue the study of rational Julia sets in \S11.
%Here are some problems for the reader.
\lln\smallskip


{\bf Concluding problems.}\smallskip

{\bf Problem 3-1.} If \[f:\hat\C\to\hat\C\] is rational of degree \[d=1\],
show that the Julia set \[J(f)\] is either vacuous, or
consists of a a single repelling or parabolic fixed point.

{\bf Problem 3-2.} If \[\Omega\subset\hat\C\]
is the basin of attraction for an attracting
periodic orbit, show that the boundary \[\partial \Omega=\bar\Omega\ssm
\Omega\] is equal to the entire Julia set. (Compare Theorem 3.3.) Here it
is essential that we include all connected components of this basin.\smallskip

{\bf Problem 3-3.} Show that a rational map \[f\] is actually a polynomial
if and only if the point at infinity is a grand orbit finite fixed
point for \[f\]. Show that \[f\] has both zero and
infinity as grand orbit finite points if and only if \[f(z)=\alpha z^n\], where
\[n\] can be any non-zero integer, negative or positive, and where
\[\alpha\ne 0\]. Conclude that
a rational map has grand orbit finite points if and only if it is conjugate,
under some fractional linear change of coordinates, either to a polynomial
or to the map \[z\mapsto 1/z^d\] for some \[d\ge 2\].\smallskip

{\bf Problem 3-4.} Using a hand calculator if necessary, decide what maps to
what in Figure 1d.\smallskip

{\bf Problem 3-5.} If
\[f(z)=z^2-6\], show that \[J(f)\] is a Cantor set contained in the intervals
\[[-3,-\sqrt 3]\cup[\sqrt 3,3]\]. More precisely, show that a point
in \[J(f)\] with\break orbit
\[z_0\mapsto z_1\mapsto\cdots\] is uniquely determined by the sequence of
signs \[z_j/|z_j|=\pm 1\].\smallskip

{\bf Problem 3-6.} For Fatou's function \[f(z)=z^2/(2+z^2)\], show that
the entire completed real axis \[\R\cup\infty\] is contained in the basin
of attraction of the origin. Show that \[J(f)\] is a Cantor set. More
precisely, given any infinite sequence of signs \[\epsilon_0\,,\,\epsilon_1\,
,\,\ldots\] show that there is one and only one point \[z=z(\epsilon_0\,,\,
\epsilon_1\,,\,\ldots)\] which satisfies the condition that \[f^{\circ n}
(z)\] is uniformly bounded away from zero, and
belongs to the half-plane \[\epsilon_n H\] for each \[n\ge 0\]. To
this end, consider the branch \[g(z)=\sqrt{2z/(1-z)}\] of \[f^{-1}\],
mapping \[\hat\C\ssm [0,1]\] onto the upper half-plane \[H\]. Starting with
any \[z_0\not\in[0,1]\], show that the successive images  \[\epsilon_0g(
\epsilon_1g(\cdots \epsilon_ng(z_0)\cdots))\] converge to the required
point \[z\in J(f)\]. Since a Cantor set cannot separate the plane, show that
the basin of attraction of the origin is equal to \[\hat\C\ssm J(f)\].
\smallskip

{\bf Problem 3-7. Self-similarity.} With rare exceptions, any shape which is observed about
one point of the Julia set will be observed infinitely often, throughout the
Julia set. More precisely, for two points \[z\] and \[z'\] of \[J=J(f)\],
let us say that \[(J,z)\] is {\bit locally conformally isomorphic} to
\[(J,z')\]
if there exists a conformal isomorphism from a neighborhood \[N\] of \[z\]
onto a neighborhood \[N'\] of \[z'\] which carries \[z\] to \[z'\] and
\[J\cap N\] onto \[J\cap N'\]. For all but finitely many \[z_0\in J\],
show that the set of \[z\] for which \[(J,z)\] is locally conformally
isomorphic to \[(J,z_0)\] is everywhere dense in \[J\].

As an example, consider the polynomial map \[f(z)=z^3+.48\,z+.928\,i\]
of Figure 2, and explain why the fixed point \[.8\,i\] looks different
from all other points of the Julia set. (See also Example 2 of \S5.)
How many pre-images does this point have?
\vfill

%\rightline{jwm, version of 8-1-91}
\eject


\footline={\hss\tenrm 4-\folio\hss}\pageno=1

\centerline{\bf \S4. Dynamics on Other Riemann Surfaces.}\bigskip

This section will try to say something about the theory of iterated
holomorphic
mappings on an arbitrary Riemann surface. However we will
concentrate on the easy cases, and simply
refer to the literature for the hard cases.
It turns out that there are only three
Riemann surfaces \[S\] for which the study of iterated mappings
is really difficult, namely
the sphere \[\hat\C\], the plane \[\C\], and the punctured plane
\[\C\ssm\{0\}\].\smallskip

The theory of iterated holomorphic mappings from the {\bit Riemann
sphere\/} \[\hat \C\] to itself has been outlined in \S3, and
will form the main goal of all of the subsequent
sections.\smallskip

We can distinguish two different classes of holomorphic maps from the {\bit
complex
plane\/} \[\C\] to itself. A {\bit polynomial map\/} of \[\C\] extends uniquely
over the Riemann sphere \[\hat\C\]. Hence the theory of polynomial
mappings can be subsumed as a special case of the theory of rational maps
of \[\hat\C\]. (See especially \S\S17-18.) On the other hand,
{\bit transcendental mappings\/} from \[\C\] to itself form an essential
distinct and more difficult subject of study. Such mappings have
been studied for more than sixty years by
many authors, starting with Fatou himself. Even iteration of the exponential
map \[\exp:\C\to\C\] provides a number of quite challenging problems.
(See for example Lyubich, and Rees.) A
proof that the Julia set \[J(\exp)\]
is the entire plane \[\C\] is included in Devaney [Dv1].
Further information about iterated transcendental functions may be found
for example in
Baker, Devaney [Dv2], Goldberg \& Keen, and in
Eremenko \& Lyubich.\smallskip

The study of iterated maps from
the {\bit punctured plane\/} \[\C\ssm\{0\}\] to itself is also a difficult and
interesting subject. (See for example Keen.)
\smallskip

For all of the uncountably many other Riemann surfaces, it turns out
that the possible dynamical behavior
is very restricted, and fairly easy to describe.
We must distinguish two very different cases according as the surface
is a {\bit torus\/} or is {\bit Hyperbolic\/}. First consider the case of
a torus \[{\bf T}=\C/\Lambda\], where \[\Lambda\] is a lattice in \[\C\].

{\QP{\bf 4.1. Theorem.} \it Every holomorphic map
\[f:{\bf T}\to{\bf T}\] is an affine map,\break
\[f(z)\equiv \alpha z+c \;\;(\mod\Lambda)\].
The corresponding Julia set \[J(f)\] is either the empty set or the
entire torus according as \[|\alpha|\le 1\] or \[|\alpha|>1\].\medskip}

Proofs will be given at the end of this section.
Here the possible values for the derivative \[\alpha\] are sharply restricted.
(Problem 4-1.) The dynamics of such iterated affine maps on the torus are
of some interest. (See Problem 4-3, as well as Example 3 of \S5.)
However, the possibilities are so limited
that this study cannot be considered very difficult.\smallskip

Finally, suppose that \[S\] is Hyperbolic. Then we will see that
any holomorphic self-map behaves in a rather dull manner under iteration.
In particular, the Julia set is alway vacuous.
Consider first the special case of the
unit disk. The following was proved by Denjoy
in 1926, sharpening an earlier result by Wolff.

{\QP {\bf 4.2. Theorem.} {\it Let \[f:D\to D\] be any
holomorphic map. Then either

$(1)$ \[f\] is a ``rotation'' about some fixed point \[z_0\in D\], or else

$(2)$ the successive iterates \[f^{\circ n}\] converge, uniformly on
compact subsets of \[D\], to a constant function \[z\mapsto c_0\in \bar D\].}
\medskip}

Here the notation \[f^{\circ n}\] stands for the \[n$-fold composition
\[f\circ\cdots\circ f\] mapping \[D\] into itself. Note that the
limiting value \[\lim_{n\to\infty}\,f^{\circ n}(z)\] may belong to the
boundary \[\partial D\]. Similar statements hold for maps of the upper
half-plane \[H\] to itself. As examples, if \[f(w)=2w\] or if \[f(w)=w+i\]
for \[w\in H\], then evidently \[\lim_{n\to\infty}\,f^{\circ n}(w)\] is
equal to
the boundary point \[+\infty\] for every \[w\in H\]. (Here we must measure
convergence with respect to the spherical metric on the compact set
\[\bar H\subset\hat\C\].)\smallskip

The corresponding assertion for an arbitrary Hyperbolic surface
is only slightly more complicated to state.
First some definitions. If \[f:S\to S\], then the sequence of points
\[\;z\mapsto f(z)\mapsto
f^{\circ 2}(z)\mapsto\cdots\;\] is called the {\bit orbit} of the point
\[z\in S\]. A
fixed point \[f(z_0)=z_0\] is said to be {\bit attracting} if the derivative
\[\lambda=f'(z_0)\] satisfies \[|\lambda|<1\]. (To be more precise, we
must first choose some local coordinate in a neighborhood of the fixed point,
and compute this derivative in terms of this local coordinate. We will be
careless about this, since in practice our Riemann surfaces will usually
be open subsets of \[\C\].) For such an attracting fixed point, the
{\bit basin of attraction} \[\Omega=\Omega(z_0)\]
is defined to be the set of all
points \[z\in S\] for which the orbit \[z\mapsto f(z)\mapsto f^{\circ 2}(z)
\cdots\] converges towards \[z_0\]. It is not difficult
to check that \[\Omega\] is an open subset of \[S\] and that \[z_0\in
\Omega\].

{\QP {\bf 4.3. Theorem.} {\it If \[S\] is a Hyperbolic Riemann surface,
then for every holomorphic map \[f:S\to S\] the Julia set \[J(f)\]
is vacuous. Furthermore either:

\noindent \[(a)\] every orbit converges
towards a unique attracting fixed point \[f(z_0)=z_0\],

\noindent\[(b)\] every orbit diverges to infinity with respect to the
\Poincare metric on \[S\],

\noindent \[(c)\] \[f\] is an automorphism of finite order, or

\noindent\[(d)\] \[S\] is isomorphic either to a disk \[D\], a punctured
disk \[D\ssm\{0\}\], or an annulus
$$	{\cal A}_r\;=\;\{z:1<|z|<r\}\,,	$$
and \[f\]
corresponds to an irrational rotation \[z\mapsto e^{2\pi i t}z\] with
\[t\not\in{\bf Q}\].}\medskip}

\noindent
Evidently these four possibilities are mutually exclusive.
Much later, in \S13, we will want to apply this Theorem to the case
where \[S\] in an open subset of the sphere \[\hat\C\] and \[f\] is a
rational map carrying \[S\] into itself. In this case, as in 4.2, it
is convenient to study limiting values which belong to the closure of \[S\].
We can then sharpen the statement in Case (b) as follows.

{\QP{\bf 4.4. Addendum.} {\it Suppose that \[U\] is a Hyperbolic open
subset of \[\hat\C\] and that \[f:U\to U\] extends holomorphically throughout
a neighborhood of the closure \[\bar U\].
Then in Case $(b)$ above, all orbits in \[U\]
must converge within \[\hat\C\] to a single boundary fixed point
$$	f(\hat z)=\hat z\;\in\;\partial U\,.	$$
This convergence is uniform on compact subsets of \[U\].}\medskip}

\noindent (In the special case where \[U\] is the standard unit disk, of course
we have this same result without the extra hypothesis
that \[f\] extends over the boundary.)\medskip

Now let us begin the proofs.\smallskip

{\bf Proof of Lemma 4.1.} To fix our ideas, suppose that \[{\bf T}=\C/\Lambda\]
where the lattice \[\Lambda
\subset\C\] is spanned by the two numbers \[1\] and \[\tau\], and where
\[\tau\not\in\R\]. Any holomorphic map \[f:{\bf T}\to {\bf T}\] lifts to a holomorphic
map \[F:\C\to\C\] on the universal covering surface. Note first that there
exist two lattice elements \[\lambda_1\,,\,\lambda_2\in\Lambda\] so that
$$	F(z+1)\;=\;F(z)+\lambda_1\,,\quad F(z+\tau)\;=\;F(z)+\lambda_2	$$
for every \[z\in\C\]. For example we certainly have \[F(z+1)\equiv F(z)\;\;
(\mod\Lambda)\], and the difference function \[F(z+1)-F(z)\in\Lambda\] must
be constant since \[\C\] is connected and the target space \[\Lambda\] is
discrete.
Now let \[g(z)=F(z)-\lambda_1 z\], so that \[g(z+1)=g(z)\]. Then
$$	g(z+\tau)\;=\;g(z)+(\lambda_2-\lambda_1\tau)\,.	$$
Thus \[g\] gives rise to a map from
the torus \[{\bf T}\] to the infinite cylinder
$$ 	\C/(\lambda_2-\lambda_1\tau)\Z\;\;\cong\;\;\C\ssm\{0\}\;,	$$
or from the torus
to \[\C\] itself if \[\;\lambda_2-\lambda_1\tau=0\]. Using Liouville's
Theorem or the Maximum Modulus Principle, we see easily that \[g\] must be
constant Thus \[g(z)\equiv c\], and\break \[F(z)=\lambda_1 z+c\]
as required. The computation of \[J(f)\] will be left as an
exercise. (See Problems 4-2 and 4-3 below.) \QED\smallskip

We will skip over 4.2 for the moment.
The proof of Theorem 4.3 begins as follows. If we are not
in Case (b), then some orbit \[z_0\mapsto z_1
\mapsto\cdots\] must possess an accumulation point \[\hat z\in S\]. That is, we
can find integers \[n(1)<n(2)<\cdots\] so that the sequence \[\{z_{n(i)}\}\]
converges to \[\hat z\]. Consider the sequence of maps \[f^{\circ(
n(i+1)-n(i))}\], carrying \[z_{n(i)}\] to \[z_{n(i+1)}\].
By Montel's Theorem, in the version 2.6, there exists some subsequence
which converges, uniformly on compact subsets, to a holomorphic map
\[h:S\to S\]. Evidently \[h(\hat z)=\hat z\].\smallskip

{\it First suppose that \[f\] strictly contracts the \Poincare metric.} Then
\[h\] must
also contract the \Poincare metric, hence \[h\] cannot have two distinct
fixed points. But \[f\] and \[h\] commute, since \[h\] is a limit of
iterates of \[f\]; hence \[f\] must map fixed points of \[h\] to fixed
points of \[h\]. Therefore the unique fixed point \[\hat z\] of \[h\]
must also be a fixed point of \[f\].
This fixed point is attracting, since \[f\] contracts the \Poincare metric.
{\it It follows that every orbit under \[f\] converges to \[\hat z\],
so that we are in Case $(a)$ of 4.3}. For otherwise,
if \[z'\] were the closest point which did not belong to the attractive
basin of \[\hat z\], then the \Poincare distance \[\rho(\hat z\,,\,z')\] could
not strictly decrease under the map \[f\].\smallskip

{\it Suppose, on the other hand, that \[f\] preserves the \Poincare metric.}
Still assuming that \[f\] has an orbit with a finite limit point, we will
prove the following.

{\QP\bf 4.5. Lemma. {\it
Under these hypotheses, some sequence of iterates \[f^{\circ m(i)}\]
converges, uniformly on compact subsets, to the identity map of \[S\].
It follows that \[f\] is necessarily an automorphism
of \[S\].}\medskip}

{\bf Proof.} Since \[f\] preserves the \Poincare metric, it
lifts to an isometry from the universal covering \[\tilde S\]
to itself. It follows easily that \[f\] is either an automorphism
or a covering map from \[S\] onto \[S\]. On the other hand,
we know that some sequence of iterates of \[f\] converges to a
map \[h\] which has a fixed point. We can lift \[h\] to a map \[H:
\tilde S\to\tilde S\] which has a corresponding fixed point.
Since \[H\] must also preserve
the \Poincare metric, it follows that \[H\] can only be a rotation about
this fixed point. Therefore, some sequence
of iterates of \[h\] converges to the identity map \[\iota\] of \[S\].
It then follows
easily that some sequence of iterates of \[f\] converges to \[\iota\]. Therefore
\[f\] is one-to-one. For if \[f(z)=f(z')\], then \[f^{\circ m}(z)
=f^{\circ m}(z')\], and passing to the limit we have \[\iota(z)=\iota(z')\],
or in other words \[z=z'\].
Since a one-to-one covering map is clearly a conformal isomorphism,
this proves the Lemma.\QED

To complete the proof of 4.3, we must prove the following.

{\QP{\bf 4.6. Lemma.} \it If an automorphism \[f\] of a Hyperbolic surface
\[S\] has iterates \[f^{\circ m}\] which are
arbitrarily close to the identity map
(uniformly on compact subsets), then either \[f\]
has finite order, or else \[S\] is isomorphic to \[D\] or \[D\ssm\{0\}\]
or to an annulus, and \[f\] corresponds to
an irrational rotation.\medskip}

(For this Lemma, we don't really need the hypothesis that \[S\] is Hyperbolic.
In the non-Hyperbolic case, the corresponding statement would be that
\[S\] is isomorphic to either \[\C\] or \[\C\ssm\{0\}\] or \[\hat\C\],
and that \[f\] corresponds to an irrational rotation.)\smallskip

Before proving 4.6, it
will be convenient to briefly consider the more general situation of
a map \[f:S\to S'\] between different Riemann surfaces, where \[S\cong\tilde S
/\Gamma\] and \[S'\cong\tilde S'/\Gamma'\]. Such a map \[f\] lifts to a map
\[F:\tilde S\to\tilde S'\] which is unique up to composition with elements
of the group \[\Gamma'\] of deck transformations of the target space.
As in the proof of 4.1,
to each deck transformation \[\gamma\in\Gamma\] there corresponds one and only
one deck transformation \[\gamma'\in\Gamma'\] so that the identity
$$	F(\gamma(z))\;=\;\gamma'(F(z))	$$
holds for all \[z\in S\]. In fact the correspondence \[\gamma\mapsto\gamma'\]
is a group homomorphism, which can be identified with the ``induced
homomorphism'' \[f_*:\pi_1(S)\to\pi_1(S')\] between fundamental groups.

We are interested in the special case \[S=S'\], with universal covering
\[\tilde S\] isomorphic to the unit disk \[D\]. It will be convenient to
choose some compact disk \[K\subset S\] and a corresponding compact disk \[K^*\]
in the universal covering surface \[\tilde S\]. If \[f^{\circ m(j)}\] is
uniformly close to the identity map on \[K\], note
that the lifted map \[F^{\circ m(j)}\] may be far from the identity map on
\[K^*\]. However, there must exist some deck transformation
\[\gamma_j\] so that the composition \[F_j=\gamma_j\circ F^{\circ m(j)}\]
is uniformly close to the identity on \[K^*\].\smallskip

We will also need the following.

{\QP{\bf 4.7. Lemma.} {\it Let \[\Gamma\subset G\] be any discrete
subgroup of a topological group \[G\]. Then for each \[\gamma\in\Gamma\]
there exist a neighborhood \[N\] of the identity in
\[G\] so that elements \[g\in N\] satisfy \[\;g\circ\gamma\circ
g^{-1}\in\Gamma\;\] only if they commute with \[\gamma\].}\medskip}

The proof is straightforward.\QED

{\bf Proof of 4.6.} By the discussion above, we have a sequence of
automorphisms\break
\[\;F_j\,=\,\gamma_j\circ F^{\circ m(j)}\,\in\, G(\tilde S)\;\]
which converge locally uniformly to
the identity automorphism. Furthermore, for each fixed \[\gamma_0\in\Gamma\]
and each \[F_j\]
there is a corresponding \[\gamma'_j\] so that
$$	F_j\circ\gamma_0\;=\gamma'_j\circ F_j\,.	$$
By 4.7, it follows that \[F_j\] commutes with \[\gamma_0\] for large \[j\].

Now let us identify \[G(\tilde S)\] with the automorphism group \[G(D)\],
which operates not only on the open disk \[D\] but also
on the closed disk \[\bar D\]. By 1.15 we know that
two non-identity elements of \[G(D)\] commute if and only if they have exactly
the same fixed points in \[\bar D\]. Thus, if we fix some non-identity
\[\gamma\in\Gamma\], and if we exclude the case where some \[F_j\] is actually
equal to the identity automorphism, then we can say that
\[\gamma\] has exactly the same fixed points as
\[F_j\] for \[j\] sufficiently large. {\it This implies that all non-identity
elements of \[\Gamma\] have the same fixed points, so that \[\Gamma\] is a
commutative group.} Evidently, either \[\Gamma\] is trivial and \[S\cong D\],
or \[\Gamma\] is free cyclic and \[S\] is an annulus or punctured disk.
(Compare Problem 2-3.) Further details are straightforward, and will be left
to the reader.\QED

{\bf Proof of Addendum 4.4.}
Starting at some arbitrary point \[z_0\] in the connected open set
\[U\subset\hat\C\], choose a path \[p:[0,1]\to U\] from the point
\[z_0=p(0)\] to \[f(z_0)=p(1)\], and continue
this path inductively for all \[t\ge 0\] by setting \[p(t+1)=f(p(t))\].
By hypothesis, the orbit \[p(0)\mapsto p(1)\mapsto p(2)\mapsto\cdots\]
converges to infinity with respect to the \Poincare metric on \[U\].
Hence it must tend to the boundary of the open set \[U\], with respect to
the spherical metric. Let \[\delta\] be the diameter
of the image \[p[0,1]\]
in the \Poincare metric for \[U\]. Then each successive image \[p[n,\,n+1]\]
must also have diameter less than \[\delta\]. Since these sets converge
to the boundary of \[U\], it follows that the diameter of \[p[n,\,n+1]\] in
the spherical metric must tend to zero as \[n\to\infty\]. (\S2.4.)
Since \[f(p(t))=
p(t+1)\], this implies that every accumulation point of the path \[p(t)\]
as \[t\to \infty\] must be a fixed point of \[f\]. It is not difficult
to check that the set of all such accumulation points must be a connected
subset of the boundary \[\partial U\]. But our hypothesis that \[f\]
continues holomorphically throughout a neighborhood of the closure \[\bar U\]
guarantees that \[f\] can have only finitely many fixed points in \[\bar U\].
This proves that the path \[p(t)\] converges to just one fixed point \[f(\hat z)
=\hat z\in\partial U\]. Thus the orbit of \[z_0\] converges to \[\hat z\]
in the spherical metric. Using 2.2 and 2.4, it follows easily that every
orbit in \[U\] converges to \[\hat z\], and that this convergence is
uniform on compact subsets of \[U\].\QED

{\bf Proof of the Denjoy-Wolff Theorem 4.2.} We
now assume that \[S\] is the unit disk \[D\subset\C\],
but do not assume that \[f\] can be continued outside of the open disk.
The following argument is taken from a
lecture of Beardon, as communicated to me by Shishikura. For any \[\epsilon
>0\], let us approximate \[f\] by the map \[z\mapsto (1-\epsilon)f(z)\]
from \[D\] into a proper subset of itself.
Then it is not difficult to check that there is one and only one fixed
point \[z_\epsilon=(1-\epsilon)f(z_\epsilon)\]. If \[f\] itself has no fixed
point, then these \[z_\epsilon\] must tend to the boundary of the disk
as \[\epsilon\to 0\].
Let \[r(\epsilon)\] be the \Poincare distance of \[z_\epsilon\] from the
origin, and consider the closed neighborhood \[N_{r(\epsilon)}(z_\epsilon
\,,\,\rho_D)\], which contains the origin as a boundary point. By Pick's
Theorem 1.10, this neighborhood is necessarily carried into itself by the map
\[z \mapsto(1-\epsilon)f(z)\]. These neighborhoods are actually round disks
(although with a different center point) with respect to the Euclidean metric.
(Problem 1-8.)
After passing to a subsequence as \[\epsilon
\to 0\], we may assume that these disks \[N_{r(\epsilon)}(z_\epsilon
\,,\,\rho_D)\] tend to a limit disk \[N_0\], bounded  by a circle (known as a
``horocircle'') which is tangent to the boundary of \[D\] at a single
point \[\hat z\]. Now \[f\] must map \[N_0\] into itself, hence the entire
orbit of \[0\] under \[f\] must be contained in \[N_0\]. Therefore, if
this orbit has no limit point in the open disk \[D\], then it must
converge towards the point of tangency
\[\hat z\]. The argument now proceeds as in 4.4.\QED
\smallskip

\lln

We conclude with problems for the reader.
\medskip

{\bf Problem 4-1.} Given some torus \[\T=\C/\Lambda\] and some \[\alpha\in\C\],
show that there exists a holomorphic map \[f(z)\equiv
\alpha z+c\] from \[\T\] to itself with derivative
\[\alpha\] if and only if
\[\alpha\Lambda\subset\Lambda\]. Evidently an arbitrary rational
integer \[\alpha
\in\Z\] will satisfy this condition. Show that there exists such a map with
derivative \[\alpha\not\in\Z\]
if and only the ratio of two generators of \[\Lambda\] satisfies a quadratic
equation with integer coefficients. Such a torus is said to admit
``{\bit complex multiplications\/}''.\smallskip

{\bf Problem 4-2.} If \[|\alpha|=1\] but \[\alpha\ne 1\], show that \[f\] is
an automorphism of finite order, and in fact of order either 2, 3, 4, or 6.
Conclude that \[J(f)\] is vacuous.
Show that all four cases can occur, for suitably chosen \[\Lambda\].
\smallskip

{\bf Problem 4-3.} If \[\alpha\ne 0\], show that any equation of the form
\[f(z)=z_0\] has exactly \[|\alpha|^2\] solutions \[z\in\T\].
If \[\alpha\ne 1\], show that \[f\]
has exactly \[|\alpha-1|^2\] fixed points. (In particular, both
\[|\alpha|^2\] and \[|\alpha-1|^2\] are
necessarily integers.) More generally, if \[|\alpha|>1\]
show that the equation \[f^{\circ n}(z)=z\] has exactly \[|\alpha^n-1|^2\]
solutions in \[\T\]. Show that the periodic points of \[f\]
are everywhere dense in \[\T\] whenever \[\alpha\ne 0\].
Conclude that the Julia set \[J(f)\] is the entire torus whenever
\[|\alpha|>1\].
\smallskip

{\bf Problem 4-4.} On the other hand, show that a map from a Hyperbolic
surface to itself can have at most one periodic point
(necessarily a fixed point), unless every point
is periodic.
\vfill\bigskip
%\rightline{jwm, version of 8-28-91}


\eject

\footline={\hss\tenrm 5-\folio\hss}\pageno=1


\centerline{\bf \S5. Smooth Julia Sets.}\medskip

Most Julia sets tend to be complicated fractal subsets of \[\hat\C\].
(Compare Brohlin [Br]. We will give a very partial explanation for this fact
in \S6.3.) This section, however,
will be devoted to three exceptional Julia sets which are actually smooth
manifolds.\smallskip

{\bf Example 1: the Circle.} As discussed already in \S3,
the unit circle appears as Julia set for the mapping
\[z\mapsto z^{\pm n}\] for any \[n\ge 2\]. Other rational
maps with this same Julia set are described in Problem 5-1. Similarly,
the real axis \[\R\cup\infty\], as the image of the unit circle under
a conformal automorphism, can appear as a Julia set. (Problem 5-2.)\smallskip

{\bf Example 2: the Interval.} Following Ulam and Von Neumann,
consider the map \[f(z)= z^2-2\], which carries the closed interval
\[I=[-2,\,2]\] onto itself. We will show that the Julia set \[J(f)\] is
equal to this interval \[I\], and that every point outside of \[I\]
belongs to the attractive basin \[\Omega(\infty)\] of the point at infinity.
(For higher degree maps with this same property, see Problem 5-3.)
For \[z_0\in I\], it is easy to check
that both solutions of the equation \[f(z)=z_0\] belong to this interval \[I\].
Since \[I\] contains a repelling fixed point \[z=2\], it follows from
Theorem 3.9 that \[I\] contains the entire Julia set \[J(f)\]. On the other
hand, the basin \[\Omega(\infty)\] is a neighborhood of infinity
whose boundary \[\partial\Omega(\infty)\] is contained in \[J(f)\subset I\]
by Theorem 3.3. Hence everything outside of \[J(f)\] must belong to this
basin. Since \[f(I)\subset I\], it follows that \[J(f)=I\].

Here is a more constructive proof that \[J(f)=I\]. We make use of the
substitution\break
\[g(w)=w+w^{-1}\], which carries the unit circle in a two-to-one
manner onto \[I=[-2,\,2]\]. For \[z_0\ne I\], the equation \[g(w)
=z_0\] has two solutions, one of which lies inside the unit circle and one
of which lies outside. Hence \[g\] maps the exterior of the closed unit disk
isomorphically onto the complement \[\C\ssm I\].
Since the squaring map in the \[w$-plane is related to \[f\] by the identity
$$	g(w^2)=g(w)^2-2=f(g(z))\,,	$$
it follows easily that the orbit of \[z\] under \[f\] either remains
bounded or diverges to infinity according as \[z\] does or does not belong to
this interval. Again using 3.3, it follows that that \[J(f)=I\].\QED

{\bf Example 3: all of \[\hat\C\].} The rest of this section will describe
an example constructed by S. Latt\`es, shortly before his death in 1918. Given
any lattice \[\Lambda\subset \C\]
we can form the quotient torus \[\T=\C/\Lambda\], as in \S4.
Thus \[\T\] is a compact
Riemann surface, and is also an additive Lie group. Note that the automorphism
\[z\mapsto -z\] of this surface has just four fixed points. For example, if
\[\Lambda=\Z+\tau\Z\] is the lattice with basis \[1\] and \[\tau\],
where \[\tau\not\in\R\],
then the four fixed points are \[0,\,1/2,\,\tau/2,\] and \[(1+\tau)/2\] modulo
\[\Lambda\].

Now form a new Riemann
surface \[S\] as a quotient of \[\T\] by identifying each \[z\in T\] with
\[-z\]. Evidently \[S\] inherits the structure of a Riemann surface (although
it loses the group structure). In fact we can use \[(z-z_j)^2\] as a local
uniformizing parameter for \[S\] near each of the four fixed points \[z_j\].
Thus the natural map \[\T\to S\] is two-to-one, except at the four ramification
points. To compute the genus of \[S\], we use the following.\eject

{\QP {\bf 5.1. Riemann-Hurwitz Formula.} {\it Let \[\T\to S\] be a branched
covering map from one compact Riemann surface onto another. Then the number of
branch points, counted with multiplicity, is equal to \[\chi(S)d-
\chi(T)\], where \[\chi\] is the Euler characteristic and \[d\] is the
degree.}\medskip}

{\bf Sketch of Proof.} Choose some triangulation of \[S\] which includes
all critical
values (that is all images of ramification points) as vertices; and let
\[a_n(S)\] be the number of\break
\[n$-simplexes, so that \[\chi(S)=a_2(S)-a_1(S)
+a_0(S)\]. In general, each simplex of \[S\] lifts to \[d\] distinct
simplices in \[\T\]. However, if \[v\] is a critical value, then there
are too few pre-images of \[v\]. The number of missing pre-images is precisely
the number of ramification points over \[v\], each counted with an
appropriate multiplicity; and the conclusion follows. %\QED
\rlap{$\sqcup$}$\sqcap$

{\bf Remark.} This proof works also for Riemann surfaces with smooth
boundary. The Formula remains true for proper maps between non-compact
Riemann surfaces, as can be verified, for example, by means of
a direct limit argument.\smallskip

In our example, since \[\T\] is a torus with Euler characteristic
\[\chi(T)=0\],
and since there are exactly four simple branch points, we conclude
that \[2\chi(S)-\chi(T)=4\], or \[\chi(S)=2\]. {\it Using the standard
formula \[\chi=2-2g\], we conclude that \[S\] is a surface of
genus zero, isomorphic to the Riemann sphere.} (Remark: The projection
map from \[\T\] to the sphere \[S\] is, up to a choice of normalization,
known as the {\bit Weierstrass \[\wp$-function}.)\smallskip

Now consider the doubling map \[z\mapsto 2z\] on \[\T\]. This commutes with
multiplication by \[-1\], and hence induces a map \[f:S\to S\].
Since the doubling map has degree four, it follows that \[f\] is a rational
map of degree four. (More generally, in place of the doubling map, we could
use any linear map which carries the lattice \[\Lambda\] into itself,
as in Problem 4-1.)

{\QP {\bf 5.2. Latt\`es Theorem.} {\it The Julia set for this rational map \[f\] is
the entire sphere \[S\].}\medskip}

{\bf Proof.} Evidently the doubling map on \[\T\] has the property that
periodic points are everywhere dense. For example, if \[r\] and \[s\]
are any rational numbers with odd denominator, then \[r+s\tau\] is periodic.
These periodic orbits are all repelling, since the multiplier is a
power of two. Evidently \[f\] inherits the same property, and the conclusion
follows by \S3.3. (Alternatively, given a small open set \[U\subset S\],
it is not difficult to show that \[f^{\circ n}(U)\] is equal to the entire
sphere \[S\] whenever \[n\] is sufficiently large. Hence no sequence of
iterates of \[f\] can converge to a limit on any open set.)\QED

In order to pin down just which rational map \[f\] has these properties,
we must first label the points of \[S\]. The four branch
points on \[\T\] map to four ``ramification points''
on \[S\], which will play a special role. Let us choose a conformal
isomorphism from \[S\] onto \[\hat\C\] which maps the first three of these
points to \[\infty\,,\,0\,,\,1\] respectively. The fourth ramification point
must then map to some \[a\in\C\ssm\{0\,,\,1\}\]. In this way we construct a
projection map \[\pi:T\to\hat\C\] of degree two, which satisfies
\[\pi(-z)=\pi(z)\], and which has critical values
$$	\pi(0)=\infty\,,\;\; \pi(1/2)=0\,,\;\; \pi(\tau/2)=1\,,\;\;
	\pi((1+\tau)/2)=a\,.	 $$
Here \[a\] can be any number distinct from \[0\,,\,1\,,\,\infty\]. In
fact, given \[a\in\C\ssm\{0\,,\,1\}\], it is not difficult to show that
there is one and only one branched covering \[ \T' \to \hat\C\]
of degree two with precisely
\[\{\infty\,,\,0\,,\,1\,,\,a\}\] as ramification points. (Compare Appendix E.)
The Riemann-Hurwitz formula shows that
this branched covering space \[\T'\] is a torus, necessarily isomorphic to
\[\C/(\Z+\tau\Z)\] for some \[\tau\not\in\R\]. The unique deck transformation
which interchanges the two pre-images of any point must preserve the linear
structure, and hence be multiplication by \[-1\].

Now the doubling map on \[\T\] corresponds under \[\pi\] to a specific
rational map\break \[f_a:\hat\C\to\hat\C\], where
$$	f_a(\pi(z))=\pi(2z)\,,	$$
and where \[J(f_a)=\hat\C\] by 5.2.
A precise computation of this map \[f_a\] is described in
Problem 5-5 below.\smallskip

{\bf Remark.} Mary Rees has proved the existence
of many more rational maps with \[J(f)=\hat\C\]. (See also Herman.)
For any degree \[d\ge 2\],
let \[{\rm Rat}(d)\] be the complex manifold consisting of all rational
maps of degree \[d\]. Rees shows that there is a subset of \[{\rm Rat}(d)\]
of positive measure consisting of maps \[f\] which are ``ergodic". By
definition, this means that any measurable subset of \[\hat\C\] which
is fully invariant under \[f\] must either have full measure or measure zero.
It can be shown that any ergodic map must necessarily have \[J(f)=\hat\C\].
\smallskip

We will study these smooth Julia sets further in \S14.7.
\lln

{\bf Concluding problems.}\smallskip

{\bf Problem 5-1.} For any \[a\in D\] the map \[\phi_a(z)=(z-a)/(1-\bar
a z)\] carries the unit disk isomorphically onto itself. (Problem 1-2.)
A finite product of the form
$$	f(z)\;=\; e^{i\theta}
\phi_{a_1}(z) \phi_{a_2}(z)\cdots\phi_{a_n}(z) $$
with \[a_j\in D\] is
called a {\bit Blaschke product} of degree \[n\]. Evidently any such \[f\]
is a rational map which carries \[D\] onto \[D\]. If \[n\ge 2\], and if
one of the factors is \[\phi_0(z)=z\] so that \[f\] has a fixed
point at the origin, show that the Julia set \[J(f)
=J(1/f)\]
is the unit circle.\smallskip

{\bf Problem 5-2.} Newton's method applied to the polynomial equation\break
\[f(z)=z^2+1=0\] yields the rational map
$$	\nu(z)\;=\;z-f(z)/f'(z)\;=\; {\textstyle{1\over 2}}(z-1/z)	$$
from \[\hat\C=\C\cup\infty\] to itself.
Show that \[J(\nu)=\R\cup\infty\].\smallskip

{\bf Problem 5-3.} The {\bit monic Tchebycheff polynomials}
$$	P_1(z)=z\;,\quad P_2(z)=z^2-2\;,\quad P_3(z)=z^3-3z\;,\;\ldots	$$
can be defined inductively
by the formula \[P_{n+1}(z)+P_{n-1}(z)=zP_n(z)\].
Show that \[P_n(w+w^{-1})
=w^n+w^{-n}\], or equivalently that \[P_n(2\cos\theta)\;=\;2\cos(n\theta)\],
and show that \[P_m\circ P_n=P_{mn}\]. For \[n\ge 2\] show that
the Julia set of \[\pm P_n\] is the interval \[[-2,\,2]\].  For \[n\ge
3\] show that \[P_n\] has \[n-1\] distinct critical points but only two
critical values, namely \[\pm 2\].\smallskip

{\bf Problem 5-4.} Show that the Julia sets studied in this section
have the following extraordinary property.
(Compare Problems 14-2, 14-1.) For all but one or two of the
periodic orbits
\[z_0\mapsto z_1\mapsto\cdots\mapsto z_n=z_0\], the multiplier
\[\lambda=f'(z_1)\cdot\cdots\cdot f'(z_n)\] satisfies \[|\lambda|=d^n\] when \[J\]
is 1-dimensional, or \[|\lambda|=d^{n/2}\] when \[J=\hat\C\],
where \[d\] is the degree.\smallskip

{\bf Problem 5-5.} For the torus \[\T=\C/(\Z+\tau\Z)\] of Example 3, show that
the involution \[z\mapsto z+\,\scriptstyle{1/ 2}\] of \[\T\] corresponds
under \[\pi\] to the involution
\[w\mapsto a/w\] of \[\hat\C\], with fixed points \[w=\pm\sqrt{a}\].
Show that the rational map \[f=f_a\] has poles at \[\infty\,,\,0\,,\,1\,,
\,a\] and double zeros at \[\pm\sqrt{a}\].  Show that \[f\] has a fixed
point of multiplier \[\lambda=4\] at infinity, and conclude that
$$	f(w)\;=\;{(w^2-a)^2 \over 4w(w-1)(w-a)}\,\;.	$$
As an example, if \[a=-1\] then
$$	f(w)\;=\;{(w^2+1)^2\over 4w(w^2-1)}\,\;.	$$
Show that the correspondence \[\;\tau\mapsto a=a(\tau)\in\C\ssm\{0,1\}\;\]
satisfies the equations
$$	a(\tau+1)=1/a(\tau)\;,\quad a(-1/\tau)=1-a(\tau)\;,	$$
and also \[\;a(-\bar\tau)=\bar a(\tau)\,\]. Conclude, for example, that
\[\;a(i)=1/2\,\], and that\break \[a((1+i)/2)=-1\]. (This correspondence
\[\tau\mapsto
a(\tau)\] is an example of an ``elliptic modular function'', and provides
an explicit representation of the upper half-plane \[H\] as a universal
covering of the thrice punctured sphere \[\C\ssm\{0,1\}\]. Compare Ahlfors
[A1, pp.269-274].)\smallskip

{\bf Problem 5-6.} For each of the six critical points \[\omega\] of \[f\],
show that \[f(f(\omega))\] is the repelling fixed point at infinity.
(According to \S13.5, the fact that each critical orbit terminates
on a repelling cycle is enough to imply that \[J(f)=\hat\C\].)
% \bigskip\bigskip


\vfill
%\rightline{jwm, version of 8-28-91}
\end
